% =====================================================================
\chapitre{Réalisation et implémentation}
% =====================================================================

Ce chapitre présente la mise en œuvre de la conception : l'environnement de travail, le
choix des modèles, la construction, l'entraînement et l'évaluation de chacun d'eux,
l'organisation du trajet, puis les fonctionnalités transverses et la stratégie
d'exécution retenue.

Avant de détailler chaque fonctionnalité, il faut souligner une partie du travail qui ne se
voit pas dans l'application : l'outillage qui encadre le développement. Un script lance
l'ensemble des serveurs dans le bon ordre et un autre les arrête ; un autre enchaîne la
préparation des données, l'entraînement et l'évaluation d'un modèle ; la configuration
(adresse du backend, clés publiques) est centralisée dans un fichier unique passé à la
compilation de l'application ; les changements de base de données passent toujours par une
nouvelle migration ; et chaque modification est vérifiée par une suite de tests
automatiques. Ces éléments ne sont pas des fonctionnalités, mais ce sont eux qui nous ont
permis d'avancer vite sans casser ce qui fonctionnait déjà.

% ---------------------------------------------------------------------
\subsection{Environnement de travail et outils}

\subsubsection{Stack technique}

Le tableau~\ref{tab:stack} récapitule les technologies utilisées dans chaque partie du
projet.

\begin{table}[H]
\centering
\small
\begin{tabularx}{\textwidth}{|L{2.6cm}|L{4.6cm}|Y|}
\hline
\rowcolor{entete}\textbf{Couche} & \textbf{Technologie} & \textbf{Rôle} \\ \hline
Mobile & Flutter (Dart), Material 3 & Un seul code pour iOS et Android, interface moderne. \\ \hline
Carte & MapKit (vue native Swift) & Carte Apple Plans, recherche de lieux, itinéraires. \\ \hline
Données & Supabase : PostgreSQL, Auth, Storage, Realtime & Comptes, données partagées, fichiers et
mises à jour en direct. \\ \hline
Service d'IA & FastAPI (Python) & Routes typées, documentation générée automatiquement. \\ \hline
Inférence & MLX, \code{mlx-lm}, \code{mlx-vlm} & Exécution des modèles en 4 bits sur la puce du Mac. \\ \hline
Entraînement & PyTorch~\cite{pytorch}, Transformers, PEFT & Affinage LoRA sur la RTX 5080. \\ \hline
MLOps & MLflow & Suivi des essais, registre des modèles, alias « champion ». \\ \hline
Itinéraires & OSRM & Temps et distances réels entre les lieux. \\ \hline
\end{tabularx}
\caption{Stack technique du projet}
\label{tab:stack}
\end{table}

Dans l'application comme dans le backend, le code est organisé en couches qui ne dépendent
que de la couche inférieure (figure~\ref{fig:couches}). Côté mobile, les écrans ne parlent
qu'à des services, et chaque service a deux implémentations : la vraie (Supabase, HTTP,
géocodage natif) et une version en mémoire. Les tests utilisent la seconde, ce qui permet de
vérifier tous les écrans sans réseau. Côté backend, chaque route délègue à un module métier
qui charge et interroge son modèle.

\begin{figure}[H]
\centering
\begin{tikzpicture}[
  couche/.style={rectangle, rounded corners=4pt, fill=#1, text=white, font=\small\bfseries,
                 minimum width=6.4cm, minimum height=0.85cm, align=center},
  note/.style={font=\scriptsize, text=gristexte, align=left, anchor=west},
  fl/.style={-{Latex[length=2mm]}, line width=0.8pt, draw=gristexte},
  ]
  \node[font=\small\bfseries\color{bleu}] at (0,0.85) {Application mobile};
  \node[couche=lagon] (m1) at (0,0) {Écrans (pages, widgets)};
  \node[couche=lagon!80!black] (m2) at (0,-1.25) {Services (interfaces)};
  \node[couche=lagon!60!black] (m3) at (0,-2.5) {Supabase · HTTP · natif iOS};
  \node[couche=gristexte] (m4) at (0,-3.75) {Implémentations en mémoire (tests)};
  \draw[fl] (m1) -- (m2); \draw[fl] (m2) -- (m3);
  \draw[fl, dashed] (m2.south east) ++(-0.4,0) to[bend left=35] (m4.north east);

  \node[font=\small\bfseries\color{bleu}] at (7.6,0.85) {Backend d'IA};
  \node[couche=orange] (b1) at (7.6,0) {Routes FastAPI + vérification JWT};
  \node[couche=orange!80!black] (b2) at (7.6,-1.25) {Modules : assistant, receipts,\\ describe, trips};
  \node[couche=orange!60!black] (b3) at (7.6,-2.5) {Modèles MLX (chargement unique)};
  \node[couche=recevoir!80!black] (b4) at (7.6,-3.75) {Registre MLflow (champion)};
  \draw[fl] (b1) -- (b2); \draw[fl] (b2) -- (b3); \draw[fl] (b4) -- (b3);
\end{tikzpicture}
\caption{Organisation en couches de l'application et du backend}
\label{fig:couches}
\end{figure}

\subsubsection{Base de données et migrations}

Le schéma Supabase évolue uniquement par fichiers de migration datés, appliqués par un
script qui garde la trace de ceux déjà passés. Douze migrations ont été écrites au fil du
projet : le schéma initial, puis les photos, la suppression de compte, les demandes pour
rejoindre, le verrouillage de la devise, l'enregistrement atomique des dépenses,
l'archivage et enfin le programme des voyages. Une règle a été tenue du début à la fin : les
migrations ne contiennent que de la structure (tables, règles, fonctions), jamais de
données. La base livrée est vide, et ce sont les utilisateurs qui la remplissent.

\subsubsection{MLX et MLflow}

Le serveur MLflow tourne sur le Mac et sert à la fois l'interface de suivi et le stockage
des fichiers produits par les entraînements. Le PC d'entraînement y envoie ses résultats
par le réseau local. Chaque modèle a sa propre expérience (« Equipay » pour l'assistant,
« Equipay-receipts » pour le scanner, « Equipay-describe » pour la saisie par phrase,
« Equipay-trips » pour l'organisation des trajets) et son entrée dans le registre. Côté Mac,
les modèles sont convertis et exécutés avec MLX~\cite{mlx}, en utilisant les bibliothèques
\code{mlx-lm}~\cite{mlxlm} pour le texte et \code{mlx-vlm}~\cite{mlxvlm} pour la vision.

% ---------------------------------------------------------------------
\subsection{Choix des modèles d'intelligence artificielle}

Le projet utilise deux familles de modèles : un modèle de langage pour tout ce qui est
texte, et un modèle vision-langage pour les tickets. Les deux choix ont été guidés par la
même contrainte : un modèle assez petit pour être affiné sur une carte de 16 Go et servi
en local, mais assez bon en français pour comprendre des phrases familières.

\subsubsection{Modèle de langage : Qwen3-4B-Instruct-2507}

Pour l'assistant, nous avons d'abord comparé plusieurs modèles ouverts de petite taille,
tous affinés avec le même jeu de données (première version) et évalués avec le même code
(tableau~\ref{tab:choix_llm}). Sans affinage, aucun ne dépassait un score de 0,30 : ils
répondaient de façon générique et inventaient des fonctionnalités. Après affinage, Qwen3-4B
s'est montré le meilleur, et la version entraînée sur les 36 couches du modèle l'a encore
amélioré. C'est ce modèle qui a été retenu pour toutes les tâches de texte.

\begin{table}[H]
\centering
\small
\begin{tabularx}{\textwidth}{|L{4.4cm}|C{1.9cm}|C{2.2cm}|Y|C{0.9cm}|}
\hline
\rowcolor{entete}\textbf{Modèle} & \textbf{Paramètres} & \textbf{Score après affinage} & \textbf{Remarques} & \textbf{Réf.} \\ \hline
Qwen3-4B-Instruct-2507, LoRA sur les 36 couches & 4 milliards & \textbf{0,822} & Modèle retenu : meilleur score, bon
niveau en français. & \cite{qwen3_4b} \\ \hline
Qwen3-4B-Instruct-2507, LoRA partiel (Mac) & 4 milliards & 0,803 & Premier essai, LoRA sur une partie des couches, entraîné sur le Mac avec MLX. & \cite{qwen3_4b} \\ \hline
Qwen2.5-1.5B-Instruct & 1,5 milliard & 0,780 & Deux fois plus rapide, mais score inférieur. & \cite{qwen25} \\ \hline
Llama-3.2-3B-Instruct & 3 milliards & 0,770 & Score le plus faible des modèles affinés. & \cite{llama32} \\ \hline
\end{tabularx}
\caption{Comparaison des modèles de langage affinés pour l'assistant (jeu de données v1)}
\label{tab:choix_llm}
\end{table}

\subsubsection{Modèle vision-langage : Qwen3-VL-4B-Instruct}

Pour les tickets, nous avons écarté l'approche classique « OCR puis analyse du texte »,
qui perd la mise en page du ticket : la position d'un montant (en bas, en gras, à côté du
mot \textit{Total}) est justement ce qui permet de le reconnaître. Qwen3-VL-4B~\cite{qwen3vl}
lit directement l'image et répond en JSON. Sans affinage, il obtenait déjà un score de 0,792
sur notre jeu de test, ce qui en faisait un bon point de départ ; sa taille (4,5 milliards de
paramètres) reste compatible avec un affinage LoRA sur notre GPU.

\subsubsection{Inférence locale en 4 bits}

Avec le scanner servi en précision complète par PyTorch, le backend occupait environ 11 Go
de mémoire, ce qui laissait trop peu de place au reste du système sur un Mac de 18 Go. Nous avons donc converti chaque modèle en 4 bits avec MLX, puis mesuré l'effet de
cette conversion sur le même jeu de test (tableau~\ref{tab:quantification}). La perte est
faible et le gain de mémoire important : c'est la version 4 bits qui est servie, la version
PyTorch restant disponible en secours.

\begin{table}[H]
\centering
\small
\begin{tabularx}{\textwidth}{|Y|C{3.3cm}|C{3.3cm}|}
\hline
\rowcolor{entete}\textbf{Mesure (scanner de tickets)} & \textbf{PyTorch, 16 bits} & \textbf{MLX, 4 bits} \\ \hline
Score global & 0,954 & 0,941 \\ \hline
Total du ticket exact & 97,2 \% & 96,3 \% \\ \hline
Taille du modèle sur disque & 8,3 Go & 2,9 Go \\ \hline
Mémoire du backend (assistant et scanner chargés) & environ 11 Go & environ 5,8 Go \\ \hline
\end{tabularx}
\caption{Effet de la quantification en 4 bits sur le scanner de tickets}
\label{tab:quantification}
\end{table}

% ---------------------------------------------------------------------
\subsection{Assistant conversationnel}

\subsubsection{Jeu de données}

L'assistant doit adopter trois comportements : répondre aux questions sur l'application,
orienter vers le support quand la réponse n'existe pas (fonction absente, problème de
compte), et refuser poliment ce qui n'a rien à voir avec \equipay. Nous avons écrit à la main
les sources du jeu de données : 196 sujets répartis en douze catégories (activités,
dépenses, soldes, compte, thème, applications concurrentes, situations concrètes, etc.),
chacun avec une réponse de référence et des mots-clés attendus. Ces sources ont ensuite été
déclinées en 2\,423 questions formulées de manières différentes. Le jeu final compte
10\,921 exemples d'entraînement, 438 de validation et 727 de test.

Les réponses de référence décrivent l'état réel de l'application : lorsqu'une fonction n'est
pas encore développée (réinitialisation du mot de passe, invitation par e-mail, etc.), la
réponse attendue oriente vers le support plutôt que d'inventer une procédure.

\subsubsection{Entraînement}

Le modèle Qwen3-4B-Instruct-2507 a été affiné par LoRA sur ses 36 couches (rang 16,
$\alpha = 320$), sur le PC, pour cinq époques au maximum avec arrêt anticipé. La perte de
validation a atteint son minimum autour de la deuxième époque (0,052) : c'est ce point qui a
été conservé. La version obtenue est la seizième du registre \code{equipay-assistant}.

\subsubsection{Évaluation}

Une réponse générée est rarement identique mot pour mot à la référence. Nous avons donc
combiné plusieurs mesures complémentaires dans un score global :

\begin{itemize}
  \item le rappel des mots-clés attendus (poids 0,30), qui vérifie que les informations
        importantes sont présentes ;
  \item la similarité sémantique (poids 0,25), mesurée par le cosinus entre les plongements
        de la réponse et de la référence, calculés avec un modèle Sentence-BERT
        multilingue~\cite{sbert} ;
  \item le ROUGE-L~\cite{rouge} (poids 0,15), qui mesure la plus longue sous-séquence commune
        avec la référence ;
  \item le F1 macro du comportement (poids 0,30) : la réponse appartient-elle à la bonne
        classe (répondre, orienter vers le support, refuser) ?
\end{itemize}

Le BLEU~\cite{bleu} et le chrF~\cite{chrf} sont également calculés et enregistrés dans MLflow
à titre de comparaison, mais n'entrent pas dans le score. Le tableau~\ref{tab:assistant}
compare le modèle de base et la version retenue sur le jeu de test.

\begin{table}[H]
\centering
\small
\begin{tabularx}{\textwidth}{|Y|C{3.2cm}|C{3.2cm}|}
\hline
\rowcolor{entete}\textbf{Métrique (jeu de test, 727 questions)} & \textbf{Qwen3-4B de base} & \textbf{Champion v16} \\ \hline
Score global & 0,276 & \textbf{0,770} \\ \hline
Similarité sémantique & 0,496 & 0,807 \\ \hline
ROUGE-L & 0,146 & 0,697 \\ \hline
Comportement correct (exactitude) & 41,7 \% & 91,7 \% \\ \hline
Questions hors sujet auxquelles il répond quand même & 97,4 \% & 14,1 \% \\ \hline
Questions de support correctement orientées & --- & 70,5 \% \\ \hline
Temps de réponse moyen sur le Mac & --- & 1,3 s \\ \hline
\end{tabularx}
\caption{Résultats de l'assistant conversationnel avant et après affinage}
\label{tab:assistant}
\end{table}

Le gain est net sur toutes les mesures. Le modèle de base répondait à presque toutes les
questions, y compris aux recettes de cuisine, et décrivait des écrans qui n'existent pas.
Après affinage, il connaît l'application et respecte son rôle. La faiblesse restante concerne
les questions de support : sur une fonction inconnue, le modèle invente encore une réponse
dans environ 30 \% des cas au lieu d'orienter vers le support. C'est la piste d'amélioration
prioritaire pour une prochaine version du jeu de données.

% ---------------------------------------------------------------------
\subsection{Scanner de tickets de caisse}

\subsubsection{Une première version instructive}

La première version du scanner a été affinée uniquement sur CORD~\cite{cord}, un jeu de
tickets indonésiens. Son score sur ce jeu était excellent, mais en conditions réelles elle a
révélé deux défauts : elle ignorait les séparateurs décimaux (un total de 11,50 € devenait
1\,150) et annonçait systématiquement des roupies indonésiennes. Le modèle avait simplement
appris les habitudes de son unique source, où les montants sont des entiers et la devise
toujours la même. Cette erreur a orienté toute la seconde version.

\subsubsection{Jeu de données de la seconde version}

Le jeu de la seconde version réunit quatre sources : CORD, SROIE~\cite{sroie} (tickets
malaisiens), WildReceipt~\cite{wildreceipt} (photos prises dans des conditions variées) et
des tickets synthétiques que nous avons générés pour la France, le Maroc, les États-Unis, le
Royaume-Uni et la Suisse, avec des montants à virgule et des devises variées. Il compte
4\,638 tickets d'entraînement, 378 de validation et 615 de test, dans sept devises. Le modèle
doit renvoyer l'enseigne, l'adresse, la date, le total tel qu'il est écrit, le total en
nombre, la devise et la liste des articles. Comme toutes les sources n'annotent pas tous les
champs, les champs absents d'une source sont exclus du calcul de la perte, pour ne pas
apprendre au modèle à les laisser vides.

\subsubsection{Entraînement}

L'affinage LoRA (rang 16, $\alpha = 32$) porte sur les couches du modèle de langage
(projections d'attention et couches denses), la partie vision restant figée : 33 millions de
paramètres entraînés sur 4,47 milliards, soit 0,74 \%. Les images sont redimensionnées à
1\,024 pixels de côté au maximum, avec un lot effectif de 8 et trois époques au plus.
L'entraînement a été interrompu par une défaillance du GPU à l'étape 1\,595 sur 1\,740 ; le
meilleur point de validation ayant été atteint à l'étape 1\,158 et sauvegardé, aucune
perte de résultat n'en a découlé.

\subsubsection{Résultats}

Les trois modèles (base, v1 et v2) ont été évalués sur le même jeu de test
(tableau~\ref{tab:scanner} et figure~\ref{fig:scores}). Le score global pondère le total
(0,35), les articles (0,15), la devise, la date, l'enseigne et la validité du JSON (0,10
chacun), le total recopié et l'adresse (0,05 chacun).

\begin{table}[H]
\centering
\small
\begin{tabularx}{\textwidth}{|Y|C{2.4cm}|C{2.4cm}|C{2.4cm}|}
\hline
\rowcolor{entete}\textbf{Métrique (615 tickets de test)} & \textbf{Base} & \textbf{v1 (CORD seul)} & \textbf{v2 (retenue)} \\ \hline
Score global & 0,792 & 0,219 & \textbf{0,954} \\ \hline
Total exact & 71,2 \% & 21,1 \% & \textbf{97,2 \%} \\ \hline
Devise & 69,7 \% & 18,6 \% & \textbf{99,4 \%} \\ \hline
Date & 86,4 \% & 0,0 \% & \textbf{96,4 \%} \\ \hline
Enseigne & 89,4 \% & 0,0 \% & \textbf{91,3 \%} \\ \hline
Articles (F1) & 73,5 \% & 16,9 \% & \textbf{88,3 \%} \\ \hline
Ticket entièrement juste & 23,6 \% & 0,0 \% & \textbf{74,5 \%} \\ \hline
Erreur « virgule ignorée » & 0,0 \% & 72,7 \% & \textbf{0,0 \%} \\ \hline
\end{tabularx}
\caption{Comparaison du modèle de base, de la v1 et de la v2 du scanner}
\label{tab:scanner}
\end{table}

La v2 améliore toutes les mesures. Le total est exact sur 100 \% des tickets en euros et en
dirhams du jeu de test. L'analyse des erreurs du modèle de base est instructive : dans
13,5 \% des cas, il prenait pour total la somme donnée en espèces par le client, une
confusion que l'affinage a fait disparaître. Deux limites demeurent : le modèle invente
parfois une date lorsqu'aucune n'est imprimée, et environ 13 \% des articles proposés
n'existent pas sur le ticket. C'est pourquoi l'application indique l'origine de chaque champ
et laisse l'utilisateur corriger avant d'enregistrer.

% ---------------------------------------------------------------------
\subsection{Saisie d'une dépense en une phrase}

\subsubsection{Jeu de données et entraînement}

Le modèle reçoit une phrase libre (« Resto avec Sara, j'ai payé 45 € hier, 60/40 »), la
liste des membres de l'activité, la date du jour et la devise par défaut, et doit renvoyer
la dépense au format JSON : montant, devise, date, payeur, participants, répartition,
titre, enseigne et articles. Le jeu de données a été généré à partir de modèles de phrases
variés, dans quatre styles d'écriture (français, langage SMS, darija et anglais), avec
des prénoms écrits en minuscules ou avec de petites fautes : 12\,000 exemples
d'entraînement, 500 de validation et 800 de test, auxquels s'ajoutent 50 phrases écrites à
la main pour un test plus proche de l'usage réel. Le modèle est un Qwen3-4B-Instruct-2507
affiné par LoRA sur ses 36 couches pendant trois époques.

\subsubsection{Post-traitement}

Le modèle ne connaît que les noms écrits dans la phrase, alors que l'application a besoin
d'identifiants de membres. Le backend relie donc chaque nom à un membre réel, en tolérant
les fautes de frappe (figure~\ref{fig:code_matcher}) : « moi » désigne l'utilisateur
connecté, un nom exact est accepté directement, et sinon le membre le plus proche est
retenu seulement s'il est assez proche et sans ambiguïté. Un nom inconnu n'est jamais
inventé. La répartition est ensuite convertie en pourcentages dont la somme fait exactement
100 (figure~\ref{fig:code_pourcentages}).

\begin{figure}[H]
\begin{lstlisting}[language=Python]
def __call__(self, name) -> str | None:
    key = norm(name)
    if not key:
        return None
    if key in ("moi", "me", "ana", "je"):
        return self.me_id
    exact = {mid for k, mid in self.keys if k == key}
    if len(exact) == 1:
        return exact.pop()
    # Faute de frappe (« Slma » → Salma) : le nom le plus proche, s'il est assez proche et sans ambiguïté
    scored = sorted(((difflib.SequenceMatcher(None, key, k).ratio(), mid) for k, mid in self.keys), reverse=True)
    if scored and scored[0][0] >= 0.75 and all(mid == scored[0][1] for r, mid in scored if r >= scored[0][0] - 0.05):
        return scored[0][1]
    return None
\end{lstlisting}
\caption{Rattachement d'un nom cité dans la phrase à un membre de l'activité}
\label{fig:code_matcher}
\end{figure}

\begin{figure}[H]
\begin{lstlisting}[language=Python]
def to_percentages(weights: dict[str, float]) -> dict[str, float] | None:
    """Poids (parts, montants, %) → pourcentages au centième dont la somme fait exactement 100."""
    weights = {k: v for k, v in weights.items() if v > 0}
    total = sum(weights.values())
    if not weights or total <= 0:
        return None
    ids = list(weights)
    cents = [round(weights[i] / total * 10000) for i in ids]
    cents[-1] += 10000 - sum(cents)
    return {i: c / 100 for i, c in zip(ids, cents)}
\end{lstlisting}
\caption{Conversion d'une répartition quelconque en pourcentages exacts}
\label{fig:code_pourcentages}
\end{figure}

\subsubsection{Résultats}

\begin{table}[H]
\centering
\small
\begin{tabularx}{\textwidth}{|Y|C{2.7cm}|C{2.7cm}|C{2.7cm}|}
\hline
\rowcolor{entete}\textbf{Métrique (version MLX 4 bits)} & \textbf{Test (800)} & \textbf{Phrases manuelles (50)} & \textbf{Modèle de base} \\ \hline
Score global & \textbf{0,976} & \textbf{0,954} & 0,656 / 0,585 \\ \hline
Dépense entièrement juste & 72,5 \% & 66 \% & --- \\ \hline
Montant & 99,3 \% & 100 \% & --- \\ \hline
Participants (F1) & 0,994 & 0,959 & --- \\ \hline
Payeur & 97,2 \% & 97,8 \% & --- \\ \hline
Répartition & 98,8 \% & 97,8 \% & --- \\ \hline
Champs inventés & 1,6 \% & 15,6 \% & --- \\ \hline
\end{tabularx}
\caption{Résultats du modèle de saisie par phrase}
\label{tab:phrase}
\end{table}

Le tableau~\ref{tab:phrase} montre que le montant, le payeur et la répartition sont presque
toujours justes. Les phrases écrites à la main, plus libres que celles du jeu généré, font
apparaître davantage de champs inventés (souvent une enseigne ou un titre déduits à tort).
Deux cas restent difficiles : « pour nous deux » lorsqu'une autre personne a payé, et les
prix unitaires (« 2 cafés à 3,50 ») parfois compris comme le prix de la ligne. Sur le Mac, une
phrase est traitée en 1,5 à 2,5 secondes.

\begin{figure}[H]
\centering
\begin{tikzpicture}
\begin{axis}[
  ybar, bar width=0.62cm, width=0.86\textwidth, height=6.2cm,
  ymin=0, ymax=1.18, ytick={0,0.2,0.4,0.6,0.8,1},
  yticklabel style={/pgf/number format/.cd, fixed, fixed zerofill, precision=1, use comma},
  symbolic x coords={Assistant, Scanner de tickets, Saisie par phrase},
  xtick=data, x tick label style={font=\small},
  enlarge x limits=0.25,
  legend style={at={(0.5,1.04)}, anchor=south, legend columns=-1, draw=none, font=\small},
  nodes near coords, nodes near coords style={font=\scriptsize, /pgf/number format/.cd, fixed, fixed zerofill, precision=3, use comma},
  axis line style={draw=black!40}, tick style={draw=black!40},
  ymajorgrids, grid style={draw=black!8},
  ylabel={Score global}, ylabel style={font=\small},
  ]
  \addplot[fill=gristexte!60, draw=none] coordinates {(Assistant,0.276) (Scanner de tickets,0.792) (Saisie par phrase,0.656)};
  \addplot[fill=lagon, draw=none] coordinates {(Assistant,0.770) (Scanner de tickets,0.954) (Saisie par phrase,0.976)};
  \legend{Modèle de base, Modèle affiné (champion)}
\end{axis}
\end{tikzpicture}
\caption{Score global de chaque modèle avant et après affinage, sur son jeu de test}
\label{fig:scores}
\end{figure}

% ---------------------------------------------------------------------
\subsection{Organisation du trajet}

\subsubsection{Une chaîne en quatre étapes}

Confier entièrement l'organisation d'une journée à un modèle de langage aurait été risqué :
ces modèles calculent mal les distances et peuvent inventer des temps de trajet. Nous avons
donc séparé le calcul et la rédaction. Le modèle n'intervient qu'à la fin, pour expliquer un
résultat déjà calculé :

\begin{enumerate}
  \item \textbf{Temps de trajet.} Le backend demande à OSRM~\cite{osrm}, un moteur
        d'itinéraires fondé sur les données OpenStreetMap~\cite{osm}, la matrice des temps et
        des distances entre tous les lieux, à pied, à vélo ou en voiture. Si le service ne
        répond pas, les temps sont estimés à partir de la distance à vol d'oiseau, et
        l'application le signale.
  \item \textbf{Répartition par jour.} Les lieux datés restent sur leur jour. Les lieux sans
        date sont placés, du plus apprécié au moins apprécié, sur le jour dont ils sont le plus
        proches, dans la limite du nombre de visites choisi. Les lieux qui ont plus de votes
        contre que de votes pour sont écartés et listés à part.
  \item \textbf{Ordre de visite.} Pour chaque jour, le backend cherche l'ordre qui minimise le
        temps de déplacement, en respectant les heures déjà fixées par le groupe. Jusqu'à huit
        lieux, toutes les permutations sont testées ; au-delà, une heuristique du plus proche
        voisin est améliorée par la méthode 2-opt~\cite{croes} (figure~\ref{fig:code_ordre}).
        Des heures sont ensuite proposées pour les lieux qui n'en ont pas : départ à 9 h,
        90 minutes par visite, arrondi au quart d'heure.
  \item \textbf{Rédaction.} Qwen3-4B-Instruct-2507, sans affinage, reçoit le programme
        calculé et rédige un court texte par jour. Son prompt système lui interdit de
        modifier l'ordre, les heures ou les durées, et d'inventer un lieu ou une distance.
\end{enumerate}

\begin{figure}[H]
\begin{lstlisting}[language=Python]
def best_order(indices: list[int], stops: list[Stop], d: list[list[float]]) -> list[int]:
    """Ordre de visite d'une journée qui minimise le temps de trajet (chemin ouvert)."""
    if len(indices) <= EXACT_LIMIT:
        best, best_cost = None, math.inf
        for perm in itertools.permutations(indices):
            if not respects_times(list(perm), stops):
                continue
            cost = path_cost(list(perm), d)
            if cost < best_cost - 1e-9:
                best, best_cost = list(perm), cost
        return best or list(indices)
    # Heuristique : plus proche voisin depuis chaque départ, puis améliorations 2-opt
    best, best_cost = list(indices), math.inf
    for start in indices:
        order, rest = [start], set(indices) - {start}
        while rest:
            nxt = min(rest, key=lambda j: d[order[-1]][j])
            order.append(nxt)
            rest.remove(nxt)
        improved = True
        while improved:
            improved = False
            for i in range(1, len(order) - 1):
                for k in range(i + 1, len(order)):
                    candidate = order[:i] + order[i : k + 1][::-1] + order[k + 1 :]
                    if path_cost(candidate, d) < path_cost(order, d) - 1e-9 and respects_times(candidate, stops):
                        order, improved = candidate, True
        if respects_times(order, stops) and path_cost(order, d) < best_cost:
            best, best_cost = order, path_cost(order, d)
    return best
\end{lstlisting}
\caption{Recherche du meilleur ordre de visite d'une journée (extrait simplifié)}
\label{fig:code_ordre}
\end{figure}

\subsubsection{Garanties et suivi}

L'ordre proposé n'est jamais plus long que l'ordre actuel du groupe : s'il n'y a rien à
gagner, l'application l'indique (« Votre ordre est déjà le meilleur »). Chaque organisation
demandée est enregistrée dans l'expérience MLflow « Equipay-trips » avec le nombre de lieux,
le temps gagné, le temps de réponse, l'origine des temps (OSRM ou estimation) et le texte
produit, ce qui permet
de suivre la qualité de la rédaction et de préparer un éventuel affinage. Enfin, le trajet
d'une journée peut être ouvert dans l'application Plans d'Apple~\cite{mapkit}, avec
toutes les étapes dans l'ordre et le mode de déplacement choisi.

% ---------------------------------------------------------------------
\subsection{Fonctionnalités transverses}

\subsubsection{Équilibrage des comptes}

À partir des soldes, l'application calcule la liste des remboursements à effectuer
(figure~\ref{fig:code_rembourse}). Les débiteurs et les créanciers sont triés par montant ;
le plus gros débiteur rembourse le plus gros créancier, et l'on recommence avec ce qui
reste. Pour $n$ personnes concernées, cette méthode produit au plus $n - 1$ virements.

\begin{figure}[H]
\begin{lstlisting}[language=Dart]
List<Settlement> settleUp(List<MemberSummary> members) {
  final debtors = [for (final m in members) if (m.balance < -0.005) (m, -m.balance)]
    ..sort((a, b) => b.$2.compareTo(a.$2));
  final creditors = [for (final m in members) if (m.balance > 0.005) (m, m.balance)]
    ..sort((a, b) => b.$2.compareTo(a.$2));
  final result = <Settlement>[];
  var d = 0, c = 0;
  var owe = debtors.isEmpty ? 0.0 : debtors[0].$2;
  var due = creditors.isEmpty ? 0.0 : creditors[0].$2;
  while (d < debtors.length && c < creditors.length) {
    final amount = math.min(owe, due);
    if (amount >= 0.01) result.add(Settlement(debtors[d].$1, creditors[c].$1, (amount * 100).round() / 100));
    owe -= amount;
    due -= amount;
    if (owe < 0.005 && ++d < debtors.length) owe = debtors[d].$2;
    if (due < 0.005 && ++c < creditors.length) due = creditors[c].$2;
  }
  return result;
}
\end{lstlisting}
\caption{Calcul des remboursements qui remettent tous les soldes à zéro}
\label{fig:code_rembourse}
\end{figure}

\subsubsection{Temps réel et mode hors ligne}

Les tables des dépenses, des photos et du programme sont publiées sur le canal temps réel de
Supabase. Lorsqu'un membre ajoute une dépense ou vote, les autres téléphones reçoivent un
événement et rechargent la vue concernée ; les rechargements rapprochés sont regroupés pour
éviter les appels inutiles. Sans réseau, un bandeau l'indique, et l'application se met à jour
d'elle-même au retour de la connexion.

\subsubsection{Rôles, archivage et demandes}

Rejoindre une activité passe par un code de huit caractères, puis par une demande que le
créateur accepte ou refuse. L'archivage est propre à chaque utilisateur : archiver un voyage
terminé ne le masque pas pour les autres. Supprimer une activité est réservé à son créateur ;
un autre membre peut seulement la quitter, à condition de ne pas avoir de dépenses.

% ---------------------------------------------------------------------
\subsection{Extraits de code significatifs et choix d'implémentation}

Plusieurs choix d'implémentation méritent d'être soulignés :

\begin{itemize}
  \item \textbf{La sécurité dans la base, pas dans l'application.} Les droits sont écrits en
        SQL, au plus près des données (figure~\ref{fig:code_rls}). Même une requête construite
        à la main avec un jeton valide ne peut pas lire ou modifier l'activité d'un autre
        groupe.
  \item \textbf{Un repli automatique pour le scanner.} Si le chargement de la version MLX
        échoue, le backend charge la version PyTorch et continue de répondre
        (figure~\ref{fig:code_repli}).
  \item \textbf{Les chiffres ne viennent jamais du modèle quand un calcul est possible.}
        Temps de trajet, ordre de visite, pourcentages et soldes sont calculés ; les modèles
        extraient ou rédigent, ils ne calculent pas.
  \item \textbf{Des services remplaçables.} Chaque service de l'application peut être
        remplacé par une version en mémoire, ce qui a permis d'écrire 62 tests d'interface
        exécutés sans réseau, complétés par 27 tests du backend.
\end{itemize}

\begin{figure}[H]
\begin{lstlisting}[language=SQL2]
create policy trip_items_insert on public.trip_items for insert to authenticated
  with check (
    created_by = (select auth.uid())
    and public.can_edit(activity_id)
    and public.has_program(activity_id)
  );

create policy trip_items_update on public.trip_items for update to authenticated
  using (created_by = (select auth.uid()) or public.is_owner(activity_id))
  with check (public.can_edit(activity_id));
\end{lstlisting}
\caption{Règles d'accès au programme : ajout par un membre, modification par l'auteur ou le créateur}
\label{fig:code_rls}
\end{figure}

\begin{figure}[H]
\begin{lstlisting}[language=Python]
def load(self) -> None:
    """Charge le modèle. En cas d'échec, l'erreur est gardée et l'API répond 503."""
    with self._loading:
        if self.ready:
            return
        engines = ["mlx", "torch"] if config.RECEIPT_ENGINE == "mlx" else ["torch"]
        for engine in engines:
            try:
                (self._load_mlx if engine == "mlx" else self._load_torch)()
                self.engine = engine
                self.error = None
                return
            except Exception as e:
                self.error = f"{type(e).__name__}: {e}"
                log.error("Scanner de reçus (%s) indisponible : %s", engine, self.error)
\end{lstlisting}
\caption{Chargement du scanner avec repli de MLX vers PyTorch (extrait simplifié)}
\label{fig:code_repli}
\end{figure}

% ---------------------------------------------------------------------
\subsection{Exécution locale et choix d'hébergement}

\subsubsection{Exécution sur le Mac}

Pendant le projet, le backend d'IA tourne sur le Mac de développement. Un script de
lancement démarre MLflow, attend qu'il soit prêt, démarre le backend, puis installe et lance
l'application sur l'iPhone ; un second script arrête l'ensemble. Le téléphone joint le
backend par le nom réseau du Mac plutôt que par son adresse IP, qui change d'un réseau à
l'autre. Le téléphone et le Mac doivent simplement être sur le même réseau.

\subsubsection{Comparaison avec un serveur GPU dans le cloud}

\begin{table}[H]
\centering
\small
\begin{tabularx}{\textwidth}{|L{3.0cm}|Y|Y|}
\hline
\rowcolor{entete}\textbf{Critère} & \textbf{Mac local (retenu)} & \textbf{Serveur GPU dans le cloud} \\ \hline
Coût & Nul : matériel déjà disponible & Facturation à l'heure, même sans utilisateur \\ \hline
Mise en place & Immédiate, scripts de lancement & Image, pilotes GPU, HTTPS et nom de domaine à configurer \\ \hline
Disponibilité & Seulement quand le Mac est allumé et sur le même réseau & Accessible partout, à tout moment \\ \hline
Confidentialité & Photos et phrases ne quittent pas le réseau local & Données envoyées à un hébergeur \\ \hline
Mémoire & 18 Go partagés, d'où la quantification en 4 bits & Plus de marge, modèles en 16 bits possibles \\ \hline
Cas d'usage idéal & Développement, démonstration, petit groupe & Publication de l'application au public \\ \hline
\end{tabularx}
\caption{Comparaison entre l'exécution locale et un serveur GPU dans le cloud}
\label{tab:hebergement}
\end{table}

Le tableau~\ref{tab:hebergement} résume ce choix. Pour un projet de démonstration utilisé par
quelques personnes, l'exécution locale offre un coût nul et garde les données sur le réseau
local, au prix d'une disponibilité limitée. Un passage au cloud ne demanderait aucun
changement dans l'application : il suffirait de modifier l'adresse du backend.

% ---------------------------------------------------------------------
\subsection{Conclusion du chapitre}

Ce chapitre a détaillé la réalisation d'\equipay : l'environnement de travail, le choix des
modèles, l'entraînement et l'évaluation de chacun d'eux, l'organisation du trajet, les
fonctionnalités transverses et la stratégie d'exécution retenue. Chaque modèle affiné fait
nettement mieux que son modèle de base sur un jeu de test mesuré de la même façon pour
toutes les versions, et chaque limite connue est compensée dans l'application par une
vérification humaine. Le chapitre suivant présente l'application obtenue à travers ses
principaux écrans.
